\chapter{Zaključak}
\label{ch:zakljucak}

U radu je projektovan, realizovan i izmeren kriptografski akcelerator koji na jednoj FPGA platformi objedinjuje tri algoritma savremene kriptografije: ECDH razmenu ključa nad krivom B-571, SHA-3 sa KMAC konstrukcijom za izvođenje ključeva i AES-GCM za zaštitu saobraćaja. Time je pokriven ceo hibridni tok iz poglavlja \ref{ch:osnove}, od uspostavljanja zajedničke tajne do zaštićenog paketa. Parametri sva tri algoritma birani su usklađeno, pa sistem u celini pruža sigurnosni nivo od 256 bitova, jer razmena ključa nosi približno 285, a izvođenje ključeva i šifrovanje po 256 bitova, dok oznaka od 128 bitova ostaje zasebna kategorija, jer se napad na nju izvodi jedino kroz živu vezu, pokušaj po pokušaj, pa njena dužina ukupni nivo ne spušta (tabela \ref{tab:nivoi_sigurnosti}). Sva tri jezgra su parametrizovana genericima, napisana u VHDL-u i izmerena po jedinstvenoj metodologiji, post-implementaciono i sa taktom kao jedinim ograničenjem (odeljak \ref{sec:rez_metodologija}).
\smallbreak
Merenja su dala jasnu sliku svakog jezgra. Kod ECDH akceleratora izmerena latencija potvrđuje procenu od šest naspram dva množenja po iteraciji merdevina, pa paralelno jezgro na približno istoj površini daje 2,7 puta kraće skalarno množenje, familija od petnaest izmerenih konfiguracija pokriva raspon od oko petnaest procenata do četiri petine čipa, a najbrža od njih opsluži oko 3875 razmena ključa u sekundi. SHA-3 jezgro sa jednom rundom po taktu postiže više od 400 MHz u svih sedam konfiguracija, a najviše 18,49 Gb/s kod varijante SHA3-256 sa magistralom od 64 bita. Propusnost u svim konfiguracijama određuju unos bloka i trajanje permutacije, uz jedan takt predaje bloka, a izvođenje ključa kroz KMAC okvir staje u tri desetinke mikrosekunde. AES-GCM akcelerator drži blok po taktu sa obe arhitekture AES jezgra koje šifruje brojački niz CTR režima, protočnom i višejezgarnom. Protočno jezgro samo dostiže oko 51 Gb/s i na celom putu podataka traži oko 1,4 puta manje logike od višejezgarnog, a višejezgarno je jeftiniji izbor za ciljeve do oko 20 Gb/s na nivou jezgra, odnosno do oko 13,6 Gb/s na celom putu podataka. Na paketima od 1504 bajta akcelerator zadržava 85 procenata pune vrednosti korisnog sadržaja. Gornja granica je 95 procenata, jer zaglavlja i oznaka nose ostatak paketa, a razliku do nje plaća pre svega punjenje jezgra novim brojačkim nizom.
\smallbreak
Granicu učestanosti celog puta podataka postavlja GCM sloj, zbog povratne petlje GHASH akumulatora, pa put daje blok po taktu i do oko 30 Gb/s. Paket od 1504 bajta kroz njega prolazi sa 108,40 korisnih bitova po taktu, a sistemski omotač ne košta ni učestanost ni propusnost. Dupleksni par, šifrovanje i dešifrovanje zajedno, staje u 41,4 procenta čipa u višejezgarnoj, odnosno 30,3 procenta u protočnoj varijanti, pa pored akceleratora na čipu ostaje mesta i za ostatak sistema. Sva tri jezgra zajedno, u izabranim radnim tačkama, staju u oko 84 procenta čipa, a sa BRAM varijantom protočnog jezgra u oko 72 (odeljak \ref{sec:rez_sinteza}).
\smallbreak
Ispravnost je proveravana po nivoima, od operacija u polju do celog sistema. Pored testova sa poznatim odgovorima po standardima, SHA-3 jezgro je provereno i na ploči pod PetaLinux sistemom preko namenskog upravljačkog programa, AES-GCM akcelerator u mrežnom sistemu sa generatorom saobraćaja, a ceo lanac simulacijom u kojoj dve strane uspostave ključ pa razmene zaštićene pakete (odeljak \ref{sec:rez_sinteza}).
\smallbreak
Pravci daljeg razvoja izloženi su u poglavlju \ref{ch:unapredjenja}, od namenske upravljačke logike okvira koja lanac zaokružuje na ploči, preko pomeranja granice propusnosti puta podataka i prelaska na snažniju platformu za najveće brzine, do postkvantne zamene razmene ključa, danas već i preporučenog puta, u kojoj Keccak permutacija, koju ovaj rad već ubrzava, ostaje središte i budućeg sistema.
